\chapter{Introdução}
\label{cap:introducao}

\section{Contextualização}
\label{sec:contexto}

O cenário contemporâneo da segurança cibernética é definido pela velocidade e
pela sofisticação crescentes das operações adversárias. Ataques de
\emph{ransomware}, campanhas de \emph{phishing} e a exploração de
vulnerabilidades tornaram-se mais frequentes e complexos, obrigando as
organizações a repensar continuamente suas estratégias de defesa
\cite{academia20}. A proliferação de fornecedores comerciais de vigilância e de
mercados clandestinos de credenciais válidas amplia o acesso a capacidades
ofensivas avançadas, de modo que atacantes alcançam, com regularidade, os
ambientes corporativos antes que qualquer alerta seja gerado.

Nesse contexto, a busca proativa por ameaças, ou \emph{threat hunting},
consolidou-se como uma prática que transcende a reatividade dos sistemas
tradicionais. Em vez de aguardar que um alerta seja disparado, os profissionais
de segurança adotam técnicas de investigação forense para examinar redes,
sistemas e estações de trabalho à procura de sinais de comprometimento
\cite{academia22}. Essa abordagem ativa permite identificar e neutralizar
indicadores de atividade maliciosa antes que sejam plenamente explorados pelos
atacantes, reduzindo a janela de oportunidade do adversário. A literatura
descreve essa inversão de postura como a transformação do alvo em caçador, com
ênfase no ciclo de vida da caça e no papel emergente da inteligência artificial
\cite{hillier2022turninghuntedhunterthreat}.

Um traço distintivo das ameaças modernas é a adoção sistemática de técnicas de
\emph{living-off-the-land} (LOLBins), nas quais o atacante abusa de binários,
\emph{scripts} e utilitários legítimos já presentes no sistema operacional — a
exemplo de \texttt{PowerShell}, \texttt{wmic}, \texttt{certutil},
\texttt{rundll32} e \texttt{mshta} — para executar suas ações. Por se valerem de
ferramentas confiáveis e assinadas pelo próprio fornecedor do sistema, essas
técnicas produzem telemetria que, isoladamente, é praticamente
indistinguível da atividade administrativa benigna. O resultado é uma evasão
eficaz das defesas baseadas em assinaturas e listas de bloqueio, deslocando o
ônus da detecção para a análise comportamental e contextual dos eventos.
Ambientes Windows, predominantes em redes corporativas e sustentados por uma
rica superfície de utilitários nativos, são especialmente suscetíveis a esse
padrão de abuso.

A esse desafio técnico soma-se um desafio humano e operacional. As soluções de
detecção contemporâneas, mesmo quando dotadas de modelos estatísticos acurados,
estão sujeitas à falácia da taxa-base: em ambientes nos quais a atividade
maliciosa é rara, mesmo classificadores de elevada acurácia geram um volume
desproporcional de falsos positivos, tornando-se operacionalmente ineficazes. O
analista de segurança, por sua vez, é submetido a uma sobrecarga cognitiva
permanente, dividido entre o triagem de alertas de baixa qualidade e a
investigação aprofundada de hipóteses relevantes. Essa sobrecarga compromete a
capacidade de resposta da organização e favorece a fadiga de alertas, na qual
eventos genuinamente críticos podem ser negligenciados em meio ao ruído.

Os agentes de inteligência artificial apresentam-se como uma resposta promissora
a esse cenário. Tais agentes são estruturas de \emph{software} que ampliam as
capacidades dos grandes modelos de linguagem (LLMs — \emph{Large Language
Models}), permitindo que estes raciocinem sobre evidências, interajam com
ferramentas externas e executem ações no sistema \cite{academia20}. O mecanismo
de atenção introduzido por Vaswani et al. \cite{vaswani2017attention} pavimentou
o caminho para esses modelos, e trabalhos subsequentes demonstraram seu potencial
de coordenação e tomada de decisão conjunta em arranjos multiagente
\cite{agashe2023llmcoordination}. Na segurança cibernética, propostas como o
HuntGPT \cite{ali2023huntgptintegratingmachinelearningbased} e modelos de geração
de hipóteses para a caça a ameaças \cite{10713082} ilustram a viabilidade de
aplicar LLMs à detecção proativa, conjugando explicabilidade e raciocínio
investigativo.

Essa autonomia, contudo, é uma faca de dois gumes. A mesma flexibilidade que
torna os agentes poderosos os expõe a uma classe de ataques que tem por alvo o
próprio agente, e não apenas o sistema defendido. Pesquisas recentes documentam o
envenenamento de memória e de bases de recuperação aumentada (RAG), a injeção
indireta de \emph{prompt} dissimulada em conteúdo externo e a amplificação de mau
funcionamento, na qual um erro inicial desencadeia um ciclo de ações inúteis que
degrada o desempenho do agente. Demonstrou-se, ainda, a possibilidade de
\emph{jailbreaks} infecciosos, em que uma única entrada adversarial se propaga
por uma população de agentes. Diante disso, qualquer arquitetura defensiva que
não modele um adversário igualmente capaz de atacar a camada de raciocínio é,
por construção, incompleta.

A governança em segurança cibernética emerge, portanto, como pilar para alinhar
as estratégias de defesa às necessidades do negócio e para coordenar, de forma
segura, o trabalho dos agentes \cite{academia21}. Frameworks consolidados — como
COBIT, NIST CSF, ISO 27001 e o recente ISO 42001 — precisam evoluir para
incorporar os riscos específicos introduzidos pela integração de tecnologias de
IA \cite{McIntosh2024,isaca2025leveraging}. A presente dissertação situa-se
nessa interseção: propõe a governança multiagente não como camada burocrática
acessória, mas como mecanismo técnico de resiliência, capaz de conter a
propagação de comprometimentos e de validar achados antes que se convertam em
decisões.

\section{Definição do Problema}
\label{sec:problema}

O problema de pesquisa pode ser enunciado pela tensão entre duas exigências
aparentemente conflitantes. De um lado, demanda-se um sistema de caça a ameaças
\emph{eficaz}, capaz de operar sobre a telemetria viva de um ambiente corporativo
Windows, detectar técnicas evasivas de \emph{living-off-the-land} e reduzir o
volume de falsos positivos que sobrecarrega os analistas. De outro, exige-se que
esse mesmo sistema seja \emph{resiliente}: como os agentes que o compõem são
suscetíveis a envenenamento de memória, injeção de \emph{prompt} e amplificação
de mau funcionamento, a automação introduz uma nova superfície de ataque que, se
explorada, pode subverter justamente a função de defesa que o sistema deveria
exercer.

Formula-se, assim, a questão central de pesquisa: \emph{como projetar uma
arquitetura multiagente para caça a ameaças que opere de forma eficaz sobre a
telemetria corporativa real e, simultaneamente, demonstre resiliência diante de
ataques adversariais dirigidos aos próprios agentes?} A resposta a essa questão
não se esgota na escolha de um modelo de linguagem mais capaz; ela depende,
sobretudo, da arquitetura de coordenação e dos mecanismos de governança que
regulam o fluxo de informação e a tomada de decisão entre os agentes.

\section{Lacuna de Pesquisa}
\label{sec:lacuna}

A revisão da literatura revela que, embora os agentes de IA tenham avançado de
tarefas supervisionadas e específicas para sistemas de raciocínio autônomo, as
arquiteturas atuais apresentam lacunas críticas quando confrontadas com um
adversário inteligente. Identificam-se, em particular, três lacunas que esta
dissertação se propõe a endereçar:

\begin{enumerate}
  \item \textbf{Ausência de orquestração multiagente robusta.} Faltam
  \emph{frameworks} que impeçam que o comprometimento de um único agente se
  propague sistemicamente, contendo a falha em vez de amplificá-la. A maioria das
  propostas concentra-se no desempenho de detecção e negligencia a contenção de
  comprometimentos internos.

  \item \textbf{Ausência de validação de integridade em tempo de execução.} Os
  mecanismos de memória e as bases de conhecimento dos agentes carecem de
  validação contínua contra envenenamento furtivo de dados, deixando o
  raciocínio do agente vulnerável a manipulações que persistem entre interações.

  \item \textbf{Autocorreção ineficaz frente à degradação comportamental.} Os
  agentes não dispõem de mecanismos confiáveis para detectar e reverter sua
  própria degradação, permanecendo expostos à amplificação de mau funcionamento e
  a ataques de negação de serviço sobre o ciclo de raciocínio.
\end{enumerate}

Soma-se a essas lacunas a escassez de avaliações conduzidas em condições
realistas — sobre telemetria viva, com instrumentação de produção e contra
adversários que atacam a camada de agentes — e a ausência de uma comparação
sistemática entre arquiteturas governadas e linhas de base sem governança. É
precisamente nessa interseção entre resiliência adversarial, validação de
integridade e operação realista que o MAS-Hunt se posiciona.

\section{Objetivos}
\label{sec:objetivos}

\subsection{Objetivo Geral}
\label{subsec:objetivo-geral}

Projetar, implementar e avaliar o MAS-Hunt, uma arquitetura de governança
multiagente para a busca proativa por ameaças cibernéticas em ambientes Windows,
capaz de operar sobre telemetria viva e de demonstrar resiliência mensurável a
ataques adversariais dirigidos aos agentes, em comparação a linhas de base sem
governança.

\subsection{Objetivos Específicos}
\label{subsec:objetivos-especificos}

Para alcançar o objetivo geral, estabelecem-se os seguintes objetivos
específicos:

\begin{enumerate}
  \item Revisar criticamente a literatura sobre caça a ameaças e detecção
  comportamental, agentes baseados em LLMs e sistemas multiagente, segurança de
  agentes (envenenamento de memória, injeção de \emph{prompt} e amplificação de
  mau funcionamento) e técnicas de \emph{living-off-the-land} mapeadas ao MITRE
  ATT\&CK.

  \item Especificar a arquitetura MAS-Hunt em três camadas — Conselho de
  Governança, Gerentes e Trabalhadores/Caçadores — detalhando os papéis dos
  agentes, os protocolos de mensagens tipadas e os cinco mecanismos de
  governança M1 (integridade de memória), M2 (validação cruzada), M3
  (resistência a injeção), M4 (monitoramento comportamental) e M5 (quarentena).

  \item Implementar o MAS-Hunt como um \emph{plugin} do Claude Code, materializado
  por meio de habilidades (\emph{skills}), comandos de barra (\emph{slash
  commands}) e \emph{hooks} que orquestram subagentes pelas capacidades nativas
  do agente principal.

  \item Construir um laboratório de avaliação instrumentado com Sysmon e Elastic
  Security sobre um domínio Windows, com um corpus de \emph{playbooks} de
  técnicas LOLBin (malicioso versus benigno) e um desenho adversarial que combina
  modos de injeção e níveis de severidade.

  \item Avaliar empiricamente o MAS-Hunt completo (C1-full) frente às condições de
  referência (C1-nohardening, C2-naive e C3-Elastic) quanto ao desempenho de
  detecção, à resiliência adversarial e à taxa de falsos positivos, com análise
  estatística por intervalos de confiança via \emph{bootstrap}.
\end{enumerate}

\section{Hipótese}
\label{sec:hipotese}

A hipótese central desta dissertação é a seguinte: \emph{a arquitetura MAS-Hunt,
ao incorporar mecanismos de governança multiagente (M1--M5), apresenta uma
redução estatisticamente significativa da suscetibilidade a ataques de
envenenamento de memória e de amplificação de mau funcionamento — medida pela
taxa de sucesso do ataque (\emph{Attack Success Rate}) — em comparação a uma
linha de base multiagente sem endurecimento e a um agente de passagem única,
preservando ou aprimorando o desempenho de detecção e reduzindo a taxa de falsos
positivos em relação às regras de detecção padrão do Elastic Security.}

Dessa hipótese central derivam-se duas hipóteses operacionais subordinadas: (i) a
governança contribui de forma marginal e mensurável para a resiliência, de modo
que a ablação dos mecanismos de endurecimento (C1-nohardening) eleva a taxa de
sucesso do ataque em relação ao sistema completo (C1-full); e (ii) a orquestração
multiagente governada reduz a taxa de falsos positivos em comparação tanto ao
agente de passagem única (C2-naive) quanto às regras padrão do Elastic
(C3-Elastic), sem prejuízo da revocação.

\section{Contribuições}
\label{sec:contribuicoes}

As principais contribuições desta dissertação são:

\begin{enumerate}
  \item \textbf{Uma síntese crítica} do panorama atual de ameaças e do
  estado da arte em vulnerabilidades de agentes de IA, articulando os riscos da
  camada de agentes (envenenamento de memória, injeção de \emph{prompt} e
  amplificação de mau funcionamento) às exigências operacionais da caça a
  ameaças em ambientes Windows.

  \item \textbf{A arquitetura MAS-Hunt}, uma proposta de orquestração multiagente
  em três camadas com cinco mecanismos de governança (M1--M5) e um
  \emph{framework} estruturado de análise (OBSERVAR $\rightarrow$ HIPOTETIZAR
  $\rightarrow$ EVIDÊNCIA $\rightarrow$ VALIDAÇÃO-CRUZADA $\rightarrow$
  CLASSIFICAR) que confere resistência a injeção adversarial.

  \item \textbf{Uma implementação concreta} do MAS-Hunt como \emph{plugin} do
  Claude Code — composto por habilidades, comandos de barra e \emph{hooks} que
  orquestram subagentes —, demonstrando que a governança multiagente pode ser
  realizada sobre capacidades nativas de orquestração do agente principal, e não
  apenas por meio de \emph{scripts} isolados.

  \item \textbf{Um arcabouço de avaliação reprodutível}, incluindo o laboratório
  instrumentado, o corpus de \emph{playbooks} LOLBin, o desenho adversarial e o
  pipeline de medição com \emph{health-gate} de telemetria, que permite comparar
  arquiteturas governadas e não governadas sob métricas e análise estatística
  bem definidas.
\end{enumerate}

\section{Organização do Texto}
\label{sec:organizacao}

O restante desta dissertação está organizado em sete capítulos. O
Capítulo~\ref{cap:introducao} — este capítulo — apresentou o contexto, o
problema, a lacuna de pesquisa, os objetivos, a hipótese e as contribuições do
trabalho. O Capítulo~2 desenvolve a fundamentação teórica, abrangendo a caça a
ameaças e a detecção comportamental, os LLMs e agentes, os sistemas multiagente,
a segurança de agentes, a telemetria Windows com Sysmon e Elastic Security, as
técnicas de \emph{living-off-the-land} mapeadas ao MITRE ATT\&CK e a injeção
adversarial de \emph{logs}. O Capítulo~3 examina os trabalhos relacionados,
comparando abordagens existentes e situando as lacunas que o MAS-Hunt se propõe a
preencher.

O Capítulo~4 detalha a arquitetura MAS-Hunt: a visão geral das três camadas, os
papéis dos agentes, os mecanismos M1--M5, o \emph{framework} M3 de análise, a
implementação como \emph{plugin} do Claude Code e os protocolos de mensagens
tipadas e de governança. O Capítulo~5 descreve a metodologia experimental,
incluindo o laboratório, o corpus de \emph{playbooks}, o desenho adversarial, as
quatro condições experimentais, o pipeline de medição e a análise estatística,
documentando honestamente as limitações de telemetria e o papel do
\emph{health-gate} na sua mitigação.

O Capítulo~6 apresenta os resultados obtidos; o Capítulo~7 discute suas
implicações, limitações e ameaças à validade; e o Capítulo~8 conclui o trabalho,
recapitulando as contribuições e apontando direções para pesquisas futuras.
